% Job posting: https://ca.linkedin.com/jobs/view/machine-learning-engineer-iii-data-at-acv-auctions-4436410354
\documentclass[11pt]{article}
\usepackage[left=0.7in,top=0.7in,right=0.7in,bottom=0.7in]{geometry}
\usepackage{enumitem}
\usepackage[hidelinks]{hyperref}
\usepackage{titlesec}
\usepackage{array}
\usepackage{setspace}
\usepackage{needspace}
\usepackage[T1]{fontenc}
\usepackage{newtxtext,newtxmath}
\ifdefined\pdfgentounicode
  \input{glyphtounicode}
  \pdfgentounicode=1
\fi
\setstretch{1.04}
\pagenumbering{gobble}
\titleformat{\section}{\Large\scshape}{}{4em}{}[\vspace{-0.7em}\rule{\linewidth}{0.4pt}\vspace{-0.4em}]
\titlespacing*{\section}{0pt}{1em}{0.4em}
\newenvironment{cvrole}[5]{%
  \par\noindent\textbf{\large #1} | \textit{\small #2, #3}\hfill \textit{\small #4 - #5}\par
  \begin{itemize}[leftmargin=2em, itemsep=-0.3em, topsep=-0.5em]%
}{%
  \end{itemize}\par\noindent\mbox{}\par\vspace{0.1em}%
}
\newcommand{\cvName}{Nishal Stanislaus Silva}
\newcommand{\cvEmail}{nishal.silva@hotmail.com}
\newcommand{\cvPhone}{+1 613 200 2030}
\newcommand{\cvCity}{Toronto, ON}
\newcommand{\cvSite}{https://nishal.xyz}
\newcommand{\cvLinkedIn}{https://www.linkedin.com/in/nishal-silva}
\newcommand{\cvGitHub}{https://github.com/ns2max}
\newcommand{\cvStatus}{Canadian Permanent Resident, Eligible to work in Canada without sponsorship}
\newcommand{\cvTagline}{Machine Learning Engineer | Computer Vision, Real-Time Inference \& Model Serving | PhD}
\begin{document}
\pagestyle{empty}
\begin{center}
    {\LARGE \scshape \textbf{\cvName}}{\large \textit{, PhD}}\\[1pt]
    {\small \cvTagline}\\[1pt]
    {\footnotesize\textit{\cvStatus}}\\[1pt]
    {\small
    \href{mailto:\cvEmail}{\cvEmail} \,|\, \cvPhone \,|\, \cvCity
    \,|\, \href{\cvSite}{nishal.xyz} \,|\, \href{\cvLinkedIn}{linkedin.com/in/nishal-silva}
    \,|\, \href{\cvGitHub}{github.com/ns2max}}
\end{center}

\section*{Professional Summary}
Machine learning engineer with a PhD (University of Trento, 2025) and 10+ years across industry and academia, specializing in computer vision models and their path to production. At MAS Holdings I designed and deployed OpenCV-based defect/damage-detection systems on live manufacturing lines -- multilingual care-label QC (inspection time cut 99.5\%) and loom-side fabric-defect detection (operator headcount cut 50\%) -- with C++/Python inference engines running on embedded hardware and PLC-integrated automated reject actuation. My postdoctoral research at Trento centered on exactly the problem of optimizing high-latency models for real-time inference: a published detector running at 14 ms on a Raspberry Pi 4, F1 0.76, 74x faster than the DTW baseline it replaced. I bring hands-on PyTorch, OpenCV and TensorFlow experience, backend engineering in AWS (EC2, S3, ECS), containerized services with Docker, and production ETL pipeline design (Airflow); I am extending that foundation into Kafka-based streaming specifically. Canadian Permanent Resident based in Toronto, available immediately, no sponsorship required.

\section*{Core Competencies}
\textbf{ML Engineering \& Model Deployment:} computer vision model design and training (PyTorch, OpenCV, TensorFlow), optimizing high-latency models for real-time inference, full model-lifecycle ownership (training, evaluation, deployment, monitoring), visual data annotation pipeline design, evaluation-framework design for complex visual data, containerized model serving (Docker). \\[2pt]
\textbf{Computer Vision \& Data Engineering:} industrial defect/damage detection with explainable overlays, multilingual visual QC, image registration and warping, PLC-integrated automated reject systems, high-frequency IoT and sensor ETL, backend software engineering in cloud environments (AWS EC2/S3/ECS). \\[2pt]
\textbf{Research \& Leadership:} benchmark and ablation design, quantified accuracy/latency/CPU trade-off analysis, peer-reviewed publications (JAES, IEEE), mentoring, CI/CD and code-review practice.

\section*{Technical Stack}
Python, C++17, SQL, PyTorch, OpenCV, TensorFlow, Keras, scikit-learn, AWS (EC2, S3, ECS, RDS), Docker, Airflow, ETL pipeline design, REST APIs, Git, CI/CD, pytest.

\section*{Education}
\textbf{PhD, Information Engineering and Computer Science} -- University of Trento, Italy \hfill 2021 - 2025 \\
\textbf{MSc, Telecommunication and Electronic Engineering} -- Sheffield Hallam University, UK \hfill 2016 - 2018 \\
\textbf{BEng (Hons), Electronic Engineering} -- Sheffield Hallam University, UK \hfill 2013 - 2014

\section*{Portfolio}
\textbf{Nebula} -- open-source C++17 audio-feature-extraction library with per-feature ARM/Raspberry Pi latency benchmarks (github.com/ns2max/nebula).

\section*{Experience}

\begin{cvrole}{Independent Researcher}{Self-directed}{Toronto, Canada}{Jan 2026}{Present}
\item Continued applied ML and systems work; two peer-reviewed papers accepted (Smart Drums, JAES; MIRaaS, IEEE IS2 2026), each requiring latency-optimized inference pipelines.
\item Used AI-assisted development workflows, including Claude Code, to build and maintain data and tooling pipelines (Nebula; ETL/infrastructure-as-code work on melodiq).
\end{cvrole}

\begin{cvrole}{Postdoctoral Researcher}{University of Trento}{Trento, Italy}{Jan 2025}{Dec 2025}
\item Owned the full ML pipeline for streaming audio and sensor data on the EU Horizon EIC Pathfinder project MUSMET: acquisition, feature engineering, training and evaluation, low-latency edge deployment, and monitoring.
\item Optimized a high-latency baseline model for real-time inference: published detector running at 14 ms on a Raspberry Pi 4, F1 0.76, 31.4\% CPU utilization -- 74x faster than the 1038 ms DTW baseline (JAES 2026).
\item Designed frozen-protocol benchmark suites and ablation studies quantifying accuracy/latency/CPU trade-offs, the same rigor needed to validate a vision model before production rollout.
\end{cvrole}

\begin{cvrole}{Visiting Researcher}{McGill University, IDMIL/CIRMMT}{Montreal, Canada}{May 2024}{Aug 2024}
\item Built a self-powered embedded sensor system fusing IMU and audio signals via a hierarchical finite-state machine, achieving F1 0.78 on a 500-event corpus (IEEE IS2 2025).
\end{cvrole}

\begin{cvrole}{Visiting Researcher}{University of Visual and Performing Arts}{Colombo, Sri Lanka}{Jan 2024}{Mar 2024}
\item Conducted a human-centred evaluation study of a real-time detection system, translating quantitative metrics into usability and adoption recommendations (IEEE I3DA 2025).
\end{cvrole}

\begin{cvrole}{Technical Consultant}{Forestpin (Pvt) Ltd}{Colombo, Sri Lanka}{Aug 2020}{Feb 2021}
\item Designed and deployed ML and statistical pipelines for financial forensics and risk detection on large structured enterprise datasets; built SQL + Python backend scoring services running in production.
\end{cvrole}

\begin{cvrole}{Research Engineer}{MAS Holdings (Pvt) Ltd}{Colombo, Sri Lanka}{Jan 2015}{Jul 2020}
\item Designed and trained computer vision models for automated defect/damage detection on live production floors: multilingual care-label QC using OpenCV template and feature matching with an explainable overlay, iterated from handheld scanner to industrial camera to a conveyor with PLC-controlled reject, cutting inspection time 99.5\%.
\item Built a loom-side fabric structural-defect detection system using a microscopic camera array, cutting operator headcount 50\%; owned the full model lifecycle from data collection and annotation through deployment and monitoring.
\item Ran real-time ingestion and ETL for high-frequency IoT and operational data across manufacturing sites in South Asia, North America, and Africa, backed by C++/Python inference engines on embedded hardware -- the closest prior analogue to a Kafka-style streaming pipeline in my track record, though not Kafka itself.
\item Introduced CI/CD and code-review practice; mentored engineers and delivered group-wide computer-vision training across the group.
\end{cvrole}

\section*{Selected Publications \footnotesize {\normalfont{(full list: \href{https://nishal.xyz/publications}{\textit{nishal.xyz/publications}})}}}
Silva, N. \& Turchet, L. (2026). Real-Time Audio Pattern Detection for Smart Musical Instruments. \textit{Journal of the Audio Engineering Society}, 74(3). Silva, N. \& Turchet, L. (2026, in press). Smart Drums. \textit{Journal of the Audio Engineering Society}.

\end{document}
